\documentclass[10pt,a4paper]{amsart}
\usepackage[margin=19mm]{geometry}
\usepackage{amsmath,amssymb,mathtools}
\usepackage[scr=rsfs,cal=euler]{mathalfa}
\usepackage{xfrac}
\usepackage[unicode,hidelinks,bookmarks=false]{hyperref}
\newcommand{\ie}{i.e.\ }
\newcommand{\cf}{cf.\ }
\newcommand{\resp}{respectively\ }
\newcommand{\Subsection}{\subsection}
\providecommand{\nobreakdash}{\nobreak\hbox{-}}
\setlength{\emergencystretch}{2em}
\multlinegap0pt
\usepackage[shortlabels]{enumitem}
\usepackage{xcolor}
\usepackage{algorithm}\usepackage{algpseudocode}
\newtheorem{lemma}{Lemma}
\newtheorem{proposition}{Proposition}
\title{Optimal Sensor and Actuator Selection for Factored Markov Decision Processes: Complexity, Approximability and Algorithms}
\author{Jayanth Bhargav, Mahsa Ghasemi, Shreyas Sundaram}
\date{Source: arXiv:2407.07310v3 (CC BY 4.0)}
\begin{document}
\maketitle

\noindent\emph{Source and licence.} Optimal Sensor and Actuator Selection for Factored Markov Decision Processes: Complexity, Approximability and Algorithms. Version arXiv:2407.07310v3 (CC BY 4.0), distributed by arXiv under the Creative Commons Attribution 4.0 International licence, http://creativecommons.org/licenses/by/4.0/, as recorded in the arXiv licence metadata for this exact version and shown on the abstract page https://arxiv.org/abs/2407.07310v3. Full licence notice: \texttt{licenses/arxiv-2407.07310-CC-BY-4.0.txt}.\par\smallskip

\noindent\textit{Editorial scope.} Counted unit: the article's proposition prop:1 (source0.tex lines 647--650). The article's complete original proof of it is delivered with it (source0.tex lines 651--667, one \texttt{proof} environment of 17 lines). No mathematics is written, repaired or re-derived: the statement and the proof are the article's own lines, in the article's own notation, and no retained line is substituted or re-flowed.  Retained uncounted as the source-local context the proof needs: lines 282--289.  Nothing was omitted from any retained span, so the package carries no omission record.  No retained line contains a citation, so the wrapper carries no bibliography and no bibliography entry.  No retained line contains a figure environment, an image-inclusion command or drawing code, so no image is delivered and the package declares no asset dependency.  Editorial formatting only, in addition: an AMS \texttt{amsart} preamble that restates the article's own declarations and macros used by the retained lines, namely the environments \texttt{lemma}, \texttt{proposition} on separate counters, together with the standard packages those lines need.  The retained environments print in this document as Lemma 1, Proposition 1 in order of appearance, whereas the article prints the same items with the numbering of its own full document.  The complete native source of the article as distributed in the arXiv e-print for this exact version is bundled unchanged as \texttt{sources/162089/source0.tex}.
 \begin{lemma}
 \label{lma1}
For the MDP $\Bar{\mathcal{M}}$ defined in Example 1, the following holds for any $\gamma\in (0,1)$:
\begin{itemize}
    \item [(i)] If the state of  $\Bar{\mathcal{M}}$ is measured (or observable) using the sensor $\omega$, the infinite-horizon expected return under the optimal policy is $V^*(s) = \frac{R}{(1-\gamma)}$ and the optimal policy ensures that the state of the MDP is $s_t=C$, for all $t>0$. 
    \item[(ii)] If the state of $\Bar{\mathcal{M}}$ is not measured i.e., there is no sensor installed and the agent only has access to the sequence of actions and rewards, but not the current state $s$, then the infinite-horizon expected reward beginning at belief $b_0$, under the optimal policy is $V^*(b) = 0$ and the optimal policy ensures that the state of the MDP is $s_t\notin  \{C,D\}$, for all $t>0$. 
\end{itemize}
\end{lemma} 

\begin{proposition}
\label{prop:1}
For the instance of fMDP-SS problem described in Example 3, the ratio $ r_{gre}(\Gamma)$ satisfies $\lim_{h \rightarrow \infty, c \rightarrow 0} r_{gre}(\Gamma) = 0.$
\end{proposition}

\begin{proof}
We first note that the specified reward function ensures each MDP $\mathcal{M}_i$ follows the optimal policy as in Lemma \ref{lma1} in order to obtain the optimal expected return for the fMDP $\mathcal{M}$. The overall reward of the fMDP $\mathcal{M}$ depends on the individual reward terms $\mathcal{R}_1$, $\mathcal{R}_2$ and $\mathcal{R}_4$. In the first iteration, the greedy algorithm will pick $\omega_1$, because $R_1 > R_2$ by $c$. In the second iteration, greedy would pick $\omega_2$ (because $R_2 > R_3$) and terminate due to the budget constraint. Therefore, the sensor subset selected by the greedy algorithm is $\Gamma = \{ \omega_1, \omega_2 \}$. By Lemma \ref{lma1}, it follows that the infinite-horizon expected reward of the greedy algorithm is 
\begin{equation}
\label{eq:gre}
    V_{\Gamma}^{gre} = \frac{R_1}{1-\gamma} + \frac{R_2}{1-\gamma}  = \frac{2R_2 + c}{1-\gamma}.
\end{equation}
Consider the following selection of sensors for the fMDP-SS instance:  $\Gamma = \{ \omega_2, \omega_3 \}$. By selecting sensors $\omega_2$ and $\omega_3$, the states $s_2 $ and $s_3$ can be measured. By Lemma \ref{lma1}, it follows that the states $s_2$ and $s_3$ are both in $C$ or $(0010)$ at all $t>0$, and thus we have
\begin{equation}
\label{eq:opt}
    V_{\Gamma}^{opt} = \frac{R_2}{1-\gamma} + \frac{R_4}{1-\gamma}  = \frac{R_2 + R_4}{1-\gamma}. 
\end{equation}
 Thus, we have 
\begin{equation*}
\begin{aligned}
    \lim_{h \rightarrow \infty, c \rightarrow 0} & r_{gre}(\Gamma) = \lim_{h \rightarrow \infty, c \rightarrow 0}  \frac{2R_2 + c}{R_2+R_4} =\lim_{h \rightarrow \infty}  \frac{2}{1+h} = 0 
\end{aligned}
\end{equation*} \end{proof}

\end{document}
